\section{Explicit rare-fast thermodynamic counterexample}
\label{app:rare-fast}

This appendix gives a self-contained bidirectionally connected three-state family showing that stationary thermodynamic resources do not bound the local temporal scale. Here bidirectionally connected means that each transition in the thermodynamic network has reverse-transition support; the stationary process need not satisfy detailed balance and generally has nonzero currents and entropy production.

Consider states $0,1,2$ with rates
\begin{equation}
0\xrightleftharpoons[bR]{aR}1,
\qquad
1\xrightleftharpoons[q]{cR}2,
\qquad
2\xrightleftharpoons[sR]{p}0,
\label{eq:rareRates}
\end{equation}
where $a,b,c,q,p,s>0$ are fixed and $R\to\infty$.  The optical rate ratio $a/b$ is independent of $R$, so a fixed optical detailed-balance ratio does not prevent the scaling.

Define
\begin{align}
D&=ac+bs+cs,\\
E&=ap+aq+bp+bq+cp+qs,\\
A_0&=bp+bq+cp,\\
A_1&=ap+aq+qs.
\end{align}
The exact stationary distribution is
\begin{align}
\pi_0&=\frac{A_0}{RD+E},\\
\pi_1&=\frac{A_1}{RD+E},\\
\pi_2&=\frac{RD}{RD+E}.
\end{align}
Hence
\begin{equation}
\pi_0\sim\frac{A_0}{DR},
\qquad
\pi_1\sim\frac{A_1}{DR},
\qquad
\pi_2\to1.
\end{equation}
The fast $O(R)$ rates act only on $O(R^{-1})$ stationary populations.

The stationary one-way flows have finite limits:
\begin{align}
aR\pi_0&\to\frac{aA_0}{D}, &
 bR\pi_1&\to\frac{bA_1}{D},\\
cR\pi_1&\to\frac{cA_1}{D}, &
 q\pi_2&\to q,\\
sR\pi_0&\to\frac{sA_0}{D}, &
 p\pi_2&\to p.
\end{align}
Therefore the total stationary one-way activity remains finite as $R\to\infty$.

The edge affinity logarithms are also independent of $R$:
\begin{align}
\ln\frac{aR\pi_0}{bR\pi_1}
&=\ln\frac{aA_0}{bA_1},\\
\ln\frac{cR\pi_1}{q\pi_2}
&=\ln\frac{cA_1}{qD},\\
\ln\frac{sR\pi_0}{p\pi_2}
&=\ln\frac{sA_0}{pD}.
\end{align}
The corresponding stationary edge currents remain finite because each is a difference of the finite one-way flows above.  Consequently the total stationary entropy-production rate
\begin{equation}
\Sigma_R
=\sum_{e}J_e\ln\frac{\phi_e^+}{\phi_e^-}
\end{equation}
remains bounded as $R\to\infty$.

The useful forward optical traffic is likewise finite and nonzero:
\begin{equation}
f_R=aR\pi_0
\longrightarrow\frac{aA_0}{D}>0.
\end{equation}

Now condition on an optical capture that places the detector in state $1$.  Let the transition $1\to2$ be the successful primary electrical-registration route, while $1\to0$ is a failed return.  The total holding rate in state $1$ is
\begin{equation}
\lambda_1=(b+c)R,
\end{equation}
and the probability that the first exit is the successful $1\to2$ route is
\begin{equation}
P_{\rm reg}=\frac{c}{b+c}>0,
\end{equation}
independent of $R$.  In a continuous-time Markov chain the holding time and exit destination are independent.  Therefore, conditional on successful registration,
\begin{equation}
D_R\sim\operatorname{Exp}[(b+c)R]
\end{equation}
with characteristic function
\begin{equation}
H_R(\omega)
=\frac{(b+c)R}{(b+c)R+i\omega}.
\end{equation}
For every fixed finite frequency band,
\begin{equation}
|H_R(\omega)|^2\to1
\end{equation}
uniformly as $R\to\infty$.  Thus the conditional temporal information bandwidth diverges linearly in $R$ while the optical detailed-balance ratio, useful optical traffic, stationary activity, and stationary entropy production remain bounded.

The missing resource is an absolute microscopic local rate or equivalent timing-concentration scale.  Stationary activity cannot supply it because activity weights fast rates by the vanishing stationary occupation of the fast states.